\chapter{Covariant Finite-Depth Refinement-Fiber Dynamics}
\label{ch:covariant}

\section{Finite-depth field content}

Let $\mathcal K_H$ be the detail fiber through depth $H$,
\begin{equation}
\mathcal K_H=\bigoplus_{h=1}^{H}\Ran D_h,
\qquad
D_hD_j=\delta_{hj}D_h.
\label{eq:cov-fiber}
\end{equation}
The multiplicity of level $h$ is
\begin{equation}
d_h=(p-1)p^{h-1}.
\label{eq:cov-multiplicity}
\end{equation}
The positive refinement operator is
\begin{equation}
L_H=\Lambda^2\sum_{h=1}^{H}p^{2\sigma h}D_h.
\label{eq:cov-LH}
\end{equation}
The average field $\phi$ is a real spacetime scalar.  The detail field $\chi$ is a real spacetime scalar valued in $\mathcal K_H$.  The finite-depth regulator is physical bookkeeping here.  It is not sent to infinity until ultraviolet control is supplied.

\section{Classification of the level operator}

Let
\begin{equation}
G_H=\prod_{h=1}^{H}O(d_h)
\label{eq:cov-level-group}
\end{equation}
act independently within every level.  This is the largest real orthogonal symmetry that preserves the decomposition while making no preferred choice inside a degenerate level.

\begin{theorem}[Level-orthogonal commutant]
If a self-adjoint operator $B$ on $\mathcal K_H$ commutes with every element of $G_H$, then
\begin{equation}
B=\sum_{h=1}^{H}b_hD_h
\label{eq:cov-commutant}
\end{equation}
for real coefficients $b_h$.
\end{theorem}

\begin{proof}
Each $D_h\mathcal K_H$ is an irreducible real module for its independent full orthogonal factor.  Commutation with the central sign flip on one level removes every off-diagonal map between distinct levels.  On one level, commutation with all rotations forces the restriction of $B$ to be a scalar multiple of the identity.  Summing the restrictions gives \eqref{eq:cov-commutant}.
\end{proof}

Refinement is branching rather than an invertible similarity between levels of unequal rank.  Choose orthogonal child isometries
\begin{equation}
V_{h,r}:\Ran D_h\longrightarrow\Ran D_{h+1},
\qquad r=1,\ldots,p,
\label{eq:cov-branching}
\end{equation}
such that
\begin{equation}
V_{h,r}^{*}V_{h,r}=D_h,
\qquad
V_{h,r}V_{h,r}^{*}\perp V_{h,s}V_{h,s}^{*},
\qquad
\sum_{r=1}^{p}V_{h,r}V_{h,r}^{*}=D_{h+1}.
\label{eq:cov-branching-resolution}
\end{equation}
The dimensions agree because $d_{h+1}=pd_h$.

\begin{theorem}[Branch-homogeneous quadratic operator]
Let $B=B^*$ commute with $G_H$.  If, for every $h<H$ and every child $r$,
\begin{equation}
V_{h,r}^{*}BV_{h,r}
=p^{2\sigma}D_hBD_h,
\label{eq:cov-scale-law}
\end{equation}
then
\begin{equation}
B=b_1\sum_{h=1}^{H}p^{2\sigma(h-1)}D_h.
\label{eq:cov-branch-solution}
\end{equation}
If $B\ge0$, then $b_1\ge0$.
\end{theorem}

\begin{proof}
The commutant theorem gives $B=\sum_hb_hD_h$.  Pulling the level-$h+1$ restriction back along any child isometry turns \eqref{eq:cov-scale-law} into $b_{h+1}=p^{2\sigma}b_h$.  Hence $b_h=b_1p^{2\sigma(h-1)}$.  A choice $b_1=\Lambda^2p^{2\sigma}$ reproduces \eqref{eq:cov-LH}.
\end{proof}

The numerical commutant audit used dimensions $1,2,4$.  The commutant of the independent level algebras had nullity three, one coefficient per level.  Adding the two refinement recurrences reduced the nullity to one.

\section{Leading average-detail interactions}

Assume the following EFT conditions.
\begin{enumerate}[leftmargin=2.2em]
\item The interaction is local in spacetime.
\item It is quadratic in $\chi$ and is separated into fiber orders zero and $2\sigma$.
\item It is invariant under $G_H$, while the order-$2\sigma$ part obeys branch homogeneity.
\item The average sector is invariant under $\phi\mapsto-\phi$.
\item The coefficient is analytic near $\phi=0$.
\end{enumerate}

The general term through fiber order $2\sigma$ is
\begin{equation}
-\frac12G(\phi)\ip{\chi}{\chi}_H
-\frac12F(\phi)\ip{\chi}{L_H\chi}_H,
\qquad
\begin{aligned}
G(\phi)&=m_b^2+\kappa\frac{\phi^2}{v^2}+\order(\phi^4/v^4),\\
F(\phi)&=F_0+\gamma\frac{\phi^2}{v^2}+\order(\phi^4/v^4).
\end{aligned}
\label{eq:cov-Fexpansion}
\end{equation}
After absorbing $F_0>0$ into $\Lambda$, the leading nontrivial interaction in the order-$2\sigma$ channel is
\begin{equation}
\boxed{
-\frac{\gamma}{2v^2}\phi^2\ip{\chi}{L_H\chi}_H.}
\label{eq:cov-leading-interaction}
\end{equation}

\begin{theorem}[Spectral-channel classification]
Under conditions 1--5, equation \eqref{eq:cov-leading-interaction} is the unique nonconstant interaction at quadratic order in both $\phi$ and $\chi$ within the branch-homogeneous order-$2\sigma$ channel, up to the coefficient $\gamma$.
\end{theorem}

\begin{proof}
The commutant and branch-homogeneity theorems force the order-$2\sigma$ bilinear to be proportional to $L_H$.  Average-field parity removes the linear term in $\phi$.  Analyticity leaves a constant followed by $\phi^2/v^2$.  The constant rescales $\Lambda$.  The remaining spectral-channel coefficient is $\gamma$.
\end{proof}

\begin{negativebox}[One-coupling boundary]
The symmetry assumptions do not remove the order-zero coupling $\kappa\phi^2\ip{\chi}{\chi}_H/(2v^2)$.  The cosmological action studied below makes the additional model choice $\kappa=0$.  The theorem therefore establishes one coefficient in the spectral channel, not uniqueness of the complete quadratic EFT.  It also does not calculate $\gamma$, select $p$ or $\sigma$, or derive $\Lambda$ from gravity.
\end{negativebox}

\section{Generally covariant action and field equations}

Adopting the printed model restriction $\kappa=0$, the finite-depth covariant system is
\begin{equation}
\boxed{
\begin{aligned}
S_H=\int_M\sqrt{-g}\,\dd^4x\Bigg[
&\frac{M_{\mathrm{Pl}}^2}{2}(R-2\Lambda_{\mathrm c})
-\frac12\nabla_\mu\phi\nabla^\mu\phi-U(\phi)\\
&-\frac12\ip{\nabla_\mu\chi}{\nabla^\mu\chi}_H
-\frac12\ip{\chi}{M_H^2(\phi)\chi}_H
\Bigg]+S_{\mathrm{vis}},
\end{aligned}}
\label{eq:cov-action}
\end{equation}
where
\begin{equation}
M_H^2(\phi)=m_b^2+F(\phi)L_H,
\qquad
F(\phi)=1+\gamma\frac{\phi^2}{v^2}.
\label{eq:cov-mass-op}
\end{equation}
The average potential used in the benchmark is
\begin{equation}
U(\phi)=\frac\lambda4(\phi^2-v^2)^2,
\qquad
\mu^2=\lambda v^2.
\label{eq:cov-average-potential}
\end{equation}

Variation gives
\begin{align}
\Ecal_\phi
&:=\Box\phi-U'(\phi)-\frac12F'(\phi)Y_H=0,
\label{eq:cov-phi-eom}\\
\Ecal_\chi
&:=\Box\chi-M_H^2(\phi)\chi=0,
\label{eq:cov-chi-eom}\\
Y_H
&:=\ip{\chi}{L_H\chi}_H.
\label{eq:cov-YH}
\end{align}
The sign convention uses metric signature $(-,+,+,+)$.

\section{Stress tensor and exchange identity}

The scalar-fiber stress tensor is
\begin{equation}
\begin{aligned}
T^{\mathrm{sf}}_{\mu\nu}={}&
\nabla_\mu\phi\nabla_\nu\phi
+\ip{\nabla_\mu\chi}{\nabla_\nu\chi}_H\\
&-g_{\mu\nu}\left[
\frac12(\nabla\phi)^2+U(\phi)
+\frac12\ip{\nabla\chi}{\nabla\chi}_H
+\frac12\ip{\chi}{M_H^2(\phi)\chi}_H
\right].
\end{aligned}
\end{equation}

\begin{theorem}[Covariant Ward identity]
For the finite-depth action \eqref{eq:cov-action},
\begin{equation}
\boxed{
\nabla^\mu T^{\mathrm{sf}}_{\mu\nu}
=\Ecal_\phi\nabla_\nu\phi
+\ip{\Ecal_\chi}{\nabla_\nu\chi}_H.}
\label{eq:cov-ward}
\end{equation}
The scalar-fiber stress tensor is conserved on shell.  A separately diffeomorphism-invariant visible sector adds its own on-shell conserved tensor.
\end{theorem}

\begin{proof}
Differentiate the displayed stress tensor.  The kinetic terms reduce to the Euler--Lagrange expressions after commuting covariant derivatives on spacetime scalars.  The explicit $\phi$ derivative of $M_H^2$ contributes $-\tfrac12F'Y_H\nabla_\nu\phi$.  That term combines with the average equation.  The remaining detail terms combine with \eqref{eq:cov-chi-eom}, which proves \eqref{eq:cov-ward}.
\end{proof}

Split $T=T^{(\phi)}+T^{(\chi)}$ by assigning the complete variable mass to the detail sector.  Then, on the two field equations,
\begin{align}
\nabla^\mu T^{(\chi)}_{\mu\nu}
&=-\frac12F'(\phi)Y_H\nabla_\nu\phi,
\label{eq:cov-exchange-chi}\\
\nabla^\mu T^{(\phi)}_{\mu\nu}
&=+\frac12F'(\phi)Y_H\nabla_\nu\phi.
\label{eq:cov-exchange-phi}
\end{align}
This is an internal transfer.  It is not a violation of covariant conservation.

\section{Exact FLRW exchange}

For homogeneous $\phi$ and detail modes, define
\begin{align}
\rho_\phi&=\frac12\dot\phi^2+U,
&p_\phi&=\frac12\dot\phi^2-U,\\
\rho_\chi&=\frac12\ip{\dot\chi}{\dot\chi}_H
+\frac12\ip{\chi}{M_H^2\chi}_H,
&p_\chi&=\frac12\ip{\dot\chi}{\dot\chi}_H
-\frac12\ip{\chi}{M_H^2\chi}_H.
\end{align}
The field equations imply
\begin{equation}
\boxed{
\dot\rho_\chi+3H(\rho_\chi+p_\chi)
=\frac12\dot F\,Y_H,}
\label{eq:cov-cont-chi}
\end{equation}
\begin{equation}
\boxed{
\dot\rho_\phi+3H(\rho_\phi+p_\phi)
=-\frac12\dot F\,Y_H.}
\label{eq:cov-cont-phi}
\end{equation}

At the locked origin, $\phi=0$ and $\dot F=0$.  Separate dust-like detail evolution and constant average potential are then compatible.  Once release begins, forcing $\rho_\chi\propto a^{-3}$ while holding the average potential fixed generally violates \eqref{eq:cov-cont-chi}--\eqref{eq:cov-cont-phi}.

The random finite-depth Ward check returned a total residual of zero at recorded precision.  Omitting the average exchange term left the nonzero residual
\begin{equation}
\BridgeWardNegativeResidual.
\end{equation}

\section{Tree-level stability}

\begin{theorem}[Finite-depth stability conditions]
Suppose the kinetic terms in \eqref{eq:cov-action} have their displayed signs.  On a field domain $\mathcal U$, the scalar-fiber sector has no ghost, gradient, or detail tachyon if
\begin{equation}
F(\phi)>0
\quad\text{and}\quad
m_b^2+F(\phi)\Lambda^2p^{2\sigma}>0
\qquad(\phi\in\mathcal U).
\label{eq:cov-stability}
\end{equation}
At the locked origin, local average stability additionally requires
\begin{equation}
\boxed{
M_\phi^2(0)
=-\mu^2+\frac\gamma{v^2}Y_H>0.}
\label{eq:cov-locking}
\end{equation}
\end{theorem}

\begin{proof}
The principal symbol is diagonal and canonical, so ghosts and gradient instabilities are absent.  The smallest detail eigenvalue occurs at $h=1$, which gives the second condition in \eqref{eq:cov-stability}.  Differentiating the potential twice at $\phi=0$ gives \eqref{eq:cov-locking}.
\end{proof}

For the polynomial truncation $F=1+\gamma\phi^2/v^2$, global boundedness without higher stabilizing operators requires $\gamma\ge0$ when $L_H$ has a positive eigenvalue.  If $\gamma<0$, sufficiently large $|\phi|$ makes $F$ negative and the potential is unbounded along a detail eigenvector.  This statement concerns the truncated EFT.  A positive nonpolynomial completion can change the global conclusion.

\Needspace{26\baselineskip}
\section{Polynomial Lorentzian first-order parent}

The metric action can be embedded into the real polynomial tetrad parent used by the foundational volume.  Let $e^I$ be a tetrad, let
\begin{equation}
\Sigma_I=\frac1{3!}\epsilon_{IJKL}e^J\wedge e^K\wedge e^L,
\qquad
\operatorname{vol}_e=\frac1{4!}\epsilon_{IJKL}e^I\wedge e^J\wedge e^K\wedge e^L,
\end{equation}
and introduce an auxiliary Lorentz vector $u^I$ and a fiber-valued Lorentz vector $w^I$.  Define
\begin{equation}
\boxed{
\begin{aligned}
S_{\mathrm{sf}}^{(1)}=\int_M\Bigg[
&\frac12\left(u_Iu^I+\ip{w_I}{w^I}_H\right)\operatorname{vol}_e
-u^I\dd\phi\wedge\Sigma_I\\
&-\ip{w^I}{\dd\chi}_H\wedge\Sigma_I
-V_H(\phi,\chi)\operatorname{vol}_e
\Bigg],
\end{aligned}}
\label{eq:cov-first-order}
\end{equation}
where
\begin{equation}
V_H=U(\phi)+\frac12\ip{\chi}{M_H^2(\phi)\chi}_H.
\end{equation}

\begin{theorem}[First-order equivalence]
On a nondegenerate tetrad chart, eliminating $u^I$ and $w^I$ from \eqref{eq:cov-first-order} gives the scalar-fiber matter part of \eqref{eq:cov-action}.
\end{theorem}

\begin{proof}
Variation with respect to the auxiliaries gives
\begin{equation}
u_I=e_I{}^\mu\partial_\mu\phi,
\qquad
w_I=e_I{}^\mu\partial_\mu\chi.
\end{equation}
Substitution completes the square and yields the two canonical kinetic terms with the potential $V_H$.
\end{proof}

The parent is real and polynomial when $F$ and $U$ are truncated polynomials.  The fields $\phi$ and $\chi$ are Lorentz scalars.  The auxiliaries are Lorentz vectors.  Diffeomorphisms act by Lie derivatives, and the internal level group commutes with spacetime gauge transformations.  The off-shell diffeomorphism--Lorentz algebra is therefore the direct scalar-matter extension of the existing parent algebra.  Classical BRST nilpotency follows with the ordinary scalar Lie action.  A quantum BV measure and anomaly calculation remain open.
